\section{Physics Geometry: Centerline Module}

\subsection{Files}

\begin{itemize}
  \item \texttt{src/physics/geometry/centerlines.py}
  \item \texttt{src/physics/geometry/types.py}
  \item \texttt{scripts/geometry/build\_geometry.py}
  \item \texttt{tests/test\_centerline\_geometry.py}
\end{itemize}

\subsection{Physical Significance}

The centerline is the one-dimensional skeleton of the microfluidic network.
It defines the topology, branch lengths, local tangent directions, and normal
directions used by later geometry and occupancy calculations. The centerline is
not itself a CFD boundary; it is a smooth geometric reference from which channel
regions and idealized domains can be constructed.

The current device has four labeled branches:

\[
\{\mathrm{inlet}, \mathrm{left}, \mathrm{right}, \mathrm{outlet}\}.
\]

The upper junction is the inlet bifurcation into the left and right loop
branches. The lower junction merges the two loop branches into the outlet.

\subsection{Coordinate and Unit Convention}

Raw centerline coordinates are loaded in pixel units. The module immediately
stores calibrated arc lengths using:

\[
s_\mu = s_{\mathrm{px}}\Delta x,
\qquad
\Delta x = 4~\mu\mathrm{m}/\mathrm{px}.
\]

For a branch represented by ordered points \(\mathbf{x}_i=(x_i,y_i)\), cumulative
arc length is:

\[
s_0 = 0,
\qquad
s_i = \sum_{j=1}^{i}\|\mathbf{x}_j-\mathbf{x}_{j-1}\|_2.
\]

The local tangent vector is estimated by centered finite difference at interior
points and one-sided difference at endpoints:

\[
\mathbf{t}_i =
\frac{\mathbf{x}_{i+1}-\mathbf{x}_{i-1}}
{\|\mathbf{x}_{i+1}-\mathbf{x}_{i-1}\|_2}.
\]

The normal convention is:

\[
\mathbf{n}_i = (-t_{y,i}, t_{x,i}).
\]

\subsection{Implementation Details}

The module:

\begin{itemize}
  \item inspects the centerline CSV schema;
  \item recognizes branch labels through aliases such as \texttt{channel};
  \item orders branch points using explicit order columns when available;
  \item falls back to coordinate-adjacency reconstruction when possible;
  \item standardizes branch orientation according to physical flow direction;
  \item validates loop branch lengths against the device YAML;
  \item saves CSV, JSON, NPZ, and PNG geometry artifacts.
\end{itemize}

The current centerline file contains labels in the \texttt{channel} column. When
nearest-neighbor path reconstruction fails for complex branch geometry, the CSV
row order is retained as the implicit point sequence. This is documented in the
runtime output and does not modify the source centerline.

\subsection{Outputs}

\begin{verbatim}
outputs/geometry/asymmetric_loop_h100/centerline_geometry.csv
outputs/geometry/asymmetric_loop_h100/geometry_summary.json
outputs/geometry/asymmetric_loop_h100/device_geometry.npz
outputs/geometry/asymmetric_loop_h100/centerline_geometry.png
\end{verbatim}

\begin{figure}[H]
  \centering
  \includegraphics[width=0.92\linewidth,height=0.75\textheight,keepaspectratio]{../outputs/geometry/asymmetric_loop_h100/centerline_geometry.png}
  \caption{Calibrated centerline geometry diagnostic. The plotted branches encode the physical inlet, left branch, right branch, and outlet topology.}
  \label{fig:centerline-geometry}
\end{figure}

\subsection{Execution}

\begin{verbatim}
.\.venv\Scripts\python.exe -m scripts.geometry.build_geometry \
  --experiment configs\experiments\video_2.yml \
  --overwrite
\end{verbatim}
